% =====================================================================
% DECEPTION COMPETITION — EXPERIMENTAL RESULTS (Option-R design).
% Supersedes the audit-design deception results in experimental_results.tex.
%
% NOTE: the METHODS for R differ from the audit-design deception_section.tex
% (score = (w.c)*trust with no 0.5 floor; PERSISTENT compounding credibility
% trust <- clip(trust + 0.1 - 1.5*over_promise) applied to every agent every
% round; no audit / no catch event / no threshold tau). deception_section.tex
% must be updated to these equations before this results section is dropped in.
%
% Data: 50 real-LLM episodes (full information, k=5), 14-dim mimics, in-process
% strategy benchmark, and a real-LLM transfer check (4 tiers x 5 episodes).
% =====================================================================

\subsection{Deception Competition}

\paragraph{Setup recap.}
Five sales agents pitch one resort to a mechanical buyer over $T=12$ rounds; each
round a fresh truth vector $t^{(r)}\in[0,1]^{5}$ is drawn and \emph{every agent
sees it in full} (so a claim above the truth is a deliberate over-statement, not
a guess). The buyer scores each agent by its trust-weighted appeal,
$s_i = (w\cdot c_i)\,\mathrm{tr}_i$, and books the top scorer each round
(winner reward $1/\lvert\mathcal{W}\rvert$, others $0$). Credibility is
\emph{persistent}: every agent's trust updates each round by
$\mathrm{tr}\leftarrow\mathrm{clip}(\mathrm{tr}+0.1-1.5\,d,\,0,\,1)$, where the
over-promise $d=\max(0,\,w\cdot(c-t))$ is the buyer's weighted disappointment.
Sustained over-claiming thus compounds credibility toward $0$ (and trust $0$
means score $0$), while honesty drifts it back up. This makes the game a test of
\emph{calibrated} deception: lie too much and you destroy the credibility lying
depends on.

\subsubsection{Model Behavior Characterization}
We characterize claim behavior from $50$ real-LLM episodes. The five models split
into two sharply different modes. \textbf{Llama is honest}: it over-claims on
essentially zero visible attributes and holds trust at $1.0$ throughout.
\textbf{GPT, Opus, Pro, and Grok deceive heavily}---over-claiming on $36$--$67\%$
of attributes (Opus highest at $67\%$)---and in doing so \emph{crater their own
credibility}, ending episodes at trust $\approx0$ (Table~\ref{tab:dec-r-strategy}).
Crucially, this is not skillful deception: the buyer's persistent audit means
that heavy, indiscriminate over-claiming destroys the trust multiplier that
makes any claim count.

The decisive behavioral signature is \emph{how the LLMs respond to their own
collapsing credibility}, which they observe each round (the prompt includes the
agent's trust trajectory). A strategic agent reduces over-claiming as trust falls,
to recover; the LLMs do the opposite. Table~\ref{tab:dec-r-escalate} and Fig.~\ref{fig:dec-r-escalation} bin each
model's per-round over-promise by its trust entering the round. GPT, Pro, and
Opus over-claim \emph{$4$--$8\times$ more} once their trust has cratered
($<0.3$) than when it is intact ($\geq0.6$)---i.e., they inflate hardest exactly
when inflation is futile (score $=$ claim $\times$ trust $\approx0$). They are
not modeling the trust\,$\to$\,score\,$\to$\,reward chain at all. Llama's flat
$0.000$ is the clean contrast: the honest option was available, and one model
took it.

\begin{table}[htbp]
\caption{Per-round over-promise $d=\max(0,w\cdot(c-t))$ binned by the agent's trust entering the round (50 real-LLM episodes). The four deceiving models do not pull back as credibility collapses---most over-claim \emph{more}. Cratered-trust bins have modest $n$ (11--66); read the direction, not the exact magnitude.}
\label{tab:dec-r-escalate}
\centering
\begin{tabular}{lccc}
\hline
Model & trust $\geq 0.6$ & $0.3$--$0.6$ & $<0.3$ (cratered) \\
\hline
Llama   & 0.000 & --- & --- \\
GPT     & 0.062 & 0.458 & \textbf{0.507} \\
Pro     & 0.108 & 0.462 & \textbf{0.543} \\
Opus    & 0.106 & 0.436 & \textbf{0.392} \\
Grok    & 0.079 & 0.315 & 0.172 \\
\hline
\end{tabular}
\end{table}

\begin{figure}[htbp]
\centering
\includegraphics[width=\columnwidth]{deception_r_escalation}
\caption{Per-round over-promise $d$ binned by the trust an agent carries into the round (50 real-LLM episodes). Llama never over-claims; GPT, Pro, and Opus over-claim \emph{more} as their credibility collapses---the opposite of recovery, and futile at trust~$\approx0$ (where score $=$ claim $\times$ trust $\approx0$).}
\label{fig:dec-r-escalation}
\end{figure}

\subsubsection{Mimic Fidelity}
\label{sec:dec-r-fidelity}
We distill each LLM into a small two-head network (lie head + claim regressor) on
a leakage-safe $14$-dimensional encoding (observed truth, own and opponents'
trust, round fraction, and a compact behavioral lag), and validate the
substitution at the level we use it---the distribution of emitted claims. Binning
each attribute's claims and comparing the real and mimic histograms by their
total-variation distance gives a pooled $\mathrm{TVD}=0.108$, below our
acceptance threshold of $0.15$ (Table~\ref{tab:dec-r-fidelity}). The key
inference detail is calibrated noise: the claim regressor emits a point estimate,
so we add per-attribute Gaussian noise equal to its training residual standard
deviation, which lets the mimic reproduce the \emph{spread} of each model's
claims rather than collapsing to the conditional mean (without it the pooled TVD
is $0.156$ and four of five models fail).

\begin{table}[htbp]
\caption{Distributional mimic fidelity: total-variation distance between real and mimic claim histograms (acceptance $\leq 0.15$; pooled $\mathrm{TVD}=0.108$, ACCEPT). Calibrated residual noise is what brings the assertive models under threshold.}
\label{tab:dec-r-fidelity}
\centering
\begin{tabular}{lc}
\hline
Model & TVD \\
\hline
Llama   & 0.001 \\
GPT     & 0.108 \\
Grok    & 0.121 \\
Pro     & 0.148 \\
Opus    & 0.163 \\
\hline
\multicolumn{2}{l}{Pooled $\mathrm{TVD}=0.108 \le 0.15$ (ACCEPT).} \\
\hline
\end{tabular}
\end{table}

\subsubsection{Strategy Comparison Against the Mimic Field}
\label{sec:dec-r-ladder}
Paralleling the auction baseline study, we pit hand-designed strategies against a
\emph{fixed} field of four behavioral mimics, $200$ episodes each, so outcomes
(including reward) are directly comparable. The strategies form a sophistication
ladder: \textbf{T1} honest; \textbf{T2} naive constant inflation; \textbf{T3}
trust-conditional (lie on the highest-value attribute only while credibility can
afford it); \textbf{T4} backloaded (honest early, inflate late, when credibility
damage has few remaining rounds to bite). Table~\ref{tab:dec-r-strategy} and Fig.~\ref{fig:dec-r-ladder} report
the result. Against the mimic field the ladder is clean and the gradient is
non-trivial: naive constant inflation (T2) self-destructs to half of fair share
($0.102$, trust $0$); honesty is the safe baseline ($0.231\approx$ the $0.20$
fair share); and \emph{calibrated} deception wins---T3 earns $0.414$ ($2\times$
fair share) and T4 earns $0.265$, both by lying selectively enough to keep trust
in the $0.57$--$0.63$ range. The game rewards reasoning about \emph{when} and
\emph{how much} to deceive.

The five LLMs, run as focals against the same field, all land \emph{between}
honest and naive---none reaches the calibrated-deception rung. Llama plays honest
($\approx$ fair share); GPT, Opus, Pro, and Grok deceive heavily, crater trust to
$\approx0$, and earn no more than honesty (GPT and Grok less). They occupy the
``tried to deceive but did it crudely'' band, never the skillful one.

\begin{table}[htbp]
\caption{Strategy performance against a fixed field of four behavioral mimics ($200$ episodes; fair share $=0.20$). Lie rate is over claimed attributes; win and reward are per round; final trust is end-of-episode. The math-tier ladder is well-formed (T2 $<$ T1 $<$ T4 $<$ T3); every LLM sits between honest and naive.}
\label{tab:dec-r-strategy}
\centering
\begin{tabular}{lcccc}
\hline
Strategy / model & lie\,\% & win & reward/rnd & trust \\
\hline
T2 naive-inflate     & 99 & 0.102 & 0.102 & 0.00 \\
T1 honest            & 0  & 0.252 & 0.231 & 1.00 \\
T4 backloaded        & 10 & 0.287 & 0.265 & 0.63 \\
T3 conditional       & 10 & 0.414 & \textbf{0.414} & 0.57 \\
\hline
GPT (LLM)            & 36 & 0.194 & 0.172 & 0.02 \\
Opus (LLM)           & 67 & 0.223 & 0.223 & 0.01 \\
Pro (LLM)            & 63 & 0.207 & 0.206 & 0.02 \\
Grok (LLM)           & 39 & 0.166 & 0.150 & 0.01 \\
Llama (LLM)          & 0  & 0.243 & 0.220 & 1.00 \\
\hline
\end{tabular}
\end{table}

\begin{figure}[htbp]
\centering
\includegraphics[width=\columnwidth]{deception_r_ladder}
\caption{Reward per round against the fixed mimic field, with $95\%$ bootstrap CIs (fair share $=0.20$, dotted). Calibrated deception (T3/T4) beats both honesty and naive inflation; every LLM lands between honest and naive, none reaching the skillful rung.}
\label{fig:dec-r-ladder}
\end{figure}

\subsubsection{Transfer to the Real LLMs}
\label{sec:dec-r-transfer}
The mimic-field ladder describes play against the \emph{surrogates}; whether it
describes play against the \emph{real} models is a separate question, and the one
that matters. We re-ran the four math tiers against a field of four real LLMs
(GPT, Opus, Pro, Grok; $5$ episodes each). \textbf{The ladder does not transfer}
(Table~\ref{tab:dec-r-transfer}, Fig.~\ref{fig:dec-r-transfer}). Two things break. First, every strategy earns
far less against the real LLMs ($0.00$--$0.13$) than against the mimics
($0.23$--$0.41$): the real models inflate so aggressively that they sweep nearly
all rounds among themselves, leaving almost nothing for any focal. Second, the
ordering collapses---honest play wins \emph{zero} rounds across all episodes
(versus $\approx$ fair share against mimics), and the sophistication gradient
(T3/T4 $\gg$ T2) flattens into noise: against the real models you must inflate
merely to win anything, but skillful and crude inflation no longer separate.

The lesson mirrors the auction study's transfer caution. The mimics pass
\emph{distributional} fidelity (Table~\ref{tab:dec-r-fidelity}) yet
\emph{distributional fidelity is not competitive-outcome fidelity}: winning
depends on claiming high exactly when it counts, and the mimics are competitively
softer than the real models even though their marginal claim histograms match.
Behavioral mimics are therefore faithful for \emph{characterizing} each model's
claim behavior, but \emph{not} a reliable proxy for \emph{ranking strategies}
against the real LLMs.

\begin{table}[htbp]
\caption{Transfer check: math-tier reward per round against the four \emph{real} LLMs ($5$ episodes each) vs.\ against the mimic field. The ladder does not transfer---honest play collapses to $0$ and the T3/T4 advantage flattens. The vs-real column has small $n$; the fine ordering among T2/T3/T4 is within noise, but the collapse of honesty and the across-the-board drop are robust.}
\label{tab:dec-r-transfer}
\centering
\begin{tabular}{lcc}
\hline
Strategy & vs.\ mimic field & vs.\ real LLMs \\
\hline
T1 honest      & 0.231 & \textbf{0.000} \\
T2 naive       & 0.102 & 0.117 \\
T3 conditional & 0.414 & 0.133 \\
T4 backloaded  & 0.265 & 0.083 \\
\hline
\end{tabular}
\end{table}

\begin{figure}[htbp]
\centering
\includegraphics[width=\columnwidth]{deception_r_transfer}
\caption{The mimic-field ladder does not transfer. Against the real LLMs every strategy earns far less, honest play wins zero, and the T3/T4 advantage flattens---distributional fidelity does not imply strategy-ranking fidelity.}
\label{fig:dec-r-transfer}
\end{figure}

\subsubsection{Synthesis}
The deception results form one coherent account, on real (validated) data:
\begin{enumerate}
\item \textbf{Behavior.} Under a persistent, crowd-disappointment audit, four of
the five LLMs deceive heavily and indiscriminately, destroying their own
credibility (trust $\to 0$); only Llama plays honest.
\item \textbf{No trust reasoning.} They do not respond to the collapsing
credibility they can observe---most \emph{escalate} over-claiming as trust
craters, when it is provably futile. They are not modeling the
trust/score/reward chain.
\item \textbf{Distillable but mode-limited.} The behavior is distillable into
faithful claim-distribution mimics (pooled $\mathrm{TVD}=0.108$); against that
mimic field, a simple \emph{calibrated} deception strategy (T3/T4) cleanly beats
both honesty and naive inflation, which the LLMs never discover.
\item \textbf{But it does not transfer.} The real models are far more aggressive
than their distributionally-faithful mimics, so honest play collapses against
them and the strategy ladder flattens---behavioral mimics compare claim
\emph{distributions} faithfully but cannot rank \emph{strategies} against the
real LLMs.
\end{enumerate}
The through-line with the auction is the same: each LLM carries a stable,
legible behavioral signature that small networks reproduce, and that simple
hand-designed strategies expose---together with the same methodological boundary,
that surrogates faithful enough to \emph{compare} are not faithful enough to
\emph{stand in for} the real models when the question is strategic ranking.
